\documentclass[11pt]{article}
\usepackage[margin=1in]{geometry}
\usepackage{amsmath,amssymb,amsthm}
\title{Disproof of Conjecture 00000006289}
\author{earthking11}
\date{October 6, 2026}
\newtheorem{proposition}{Proposition}
\begin{document}
\maketitle
\section*{A two-element counterexample}
Let $C_2=\{e,s\}$ with $s^2=e$.  Consider its trivial one-dimensional
representation
\[
 \rho(g)=[1]\qquad(g\in C_2).
\]
It is a representation because $[1][1]=[1]$, so
$\rho(gh)=\rho(g)\rho(h)$ for all $g,h$.

The associated character (matrix trace) is
\[
 \chi(e)=\operatorname{tr}[1]=1,
 \qquad
 \chi(s)=\operatorname{tr}[1]=1.
\]
In particular $s\ne e$ but $\chi(s)\ne0$.  Since $C_2$ is abelian, its two
conjugacy classes are $\{e\}$ and $\{s\}$.  Hence the support of $\chi$ is
both conjugacy classes, i.e. all of $C_2$, rather than the identity class.

\begin{proposition}
The trace of a finite-group representation need not vanish outside the
identity conjugacy class.
\end{proposition}
\begin{proof}
The trivial character of $C_2$ has value $1$ on the nonidentity element $s$.
\end{proof}

The identity-only support property does hold for the regular character: its
permutation matrices have trace zero away from the identity.  The filed
statement, however, does not restrict the trace to the regular character, so
the trivial character is within scope and disproves it.

The Lean project gives the complete two-element multiplication and support
tables and checks them in the kernel.
\end{document}
